\section{\uppercase{Experimental Methodology}} \label{sec4}
We evaluate the effectiveness of ranking models (RMs) for adaptive RAG configuration selection on QA tasks across three aspects: (i) RM ranking performance, (ii) RAG performance with selected configurations, and (iii) RAG performance when users select from the top-$K$ candidate configurations. This section describes the methodology for each aspect.

\subsection{Evaluation Items} \label{sec4.0}
Evaluation item (i) measures the alignment between predicted rankings and the ground-truth rankings.
Evaluation item (ii) assesses final response quality by running RAG with configurations selected under five priority-weight settings and comparing the results with baseline methods. We evaluate the top-ranked configuration (Proposed@1) and the average performance of the top five configurations (Proposed@5).
Evaluation item (iii) assesses the practical utility of presenting top-$K$ candidate configurations in realistic deployment scenarios. We simulate a workflow in which a user runs RAG with each of the top $K$ candidates and adopts the configuration with the highest empirical performance under the user's priority policy.

For each query, the ranking model first ranks candidate RAG configurations using the weights corresponding to the relevant user policy (Section~\ref{sec4.3}). The system then runs each of the top-$K$ configurations and evaluates the outputs against the ground-truth answers to calculate precision, recall, $F_1$ score, and execution latency. We select the best configuration among the $K$ executed candidates according to the operational priority and report the selected configuration's average performance across all evaluation queries. Under the ``$F_1$-oriented'' policy, for example, we select the candidate with the highest empirical $F_1$ score from the top-$K$ set ranked using $W_{\mathrm{f_1}}=1.0$. The ``Precision-oriented,'' ``Recall-oriented,'' and ``Time-oriented'' policies select the candidates with the highest empirical values for their respective metrics from sets ranked using $W_{\mathrm{pre}}=1.0$, $W_{\mathrm{rec}}=1.0$, and $W_{\mathrm{time}}=1.0$, respectively.

\subsection{Evaluation Datasets} \label{sec4.1}
We use four standard QA datasets: SQuAD~\cite{rajpurkar16}, HotpotQA~\cite{yang18}, 2WikiMultihopQA~\cite{ho20}, and MuSiQue~\cite{trivedi22}. Each dataset provides queries, associated documents, and ground-truth answers. We combine the datasets to evaluate performance across different difficulty levels. From each dataset, we randomly sampled 100 instances for training and 20 for evaluation, yielding 400 training instances and 80 evaluation instances in total.

\subsection{Ranking Models} \label{sec4.2}
We use LightGBM~\cite{KeGuolin17}, XGBoost~\cite{ChenTianqi16}, and CatBoost~\cite{Prokhorenkova18} as ranking algorithms. LightGBM and XGBoost capture nonlinear interactions among features, whereas CatBoost is designed to handle categorical features. For evaluation items (ii) and (iii), we use the model with the best performance on evaluation item (i).

\subsection{Proposed Method Configurations} \label{sec4.3}
For evaluation items (ii) and (iii), we use five configurations of the proposed method, each with different priority weights in Eq.~\eqref{eq1}:
1. Balanced ($F_1$-Time): $W_{\mathrm{f_1}} = W_{\mathrm{time}} = 0.5$.
2. Precision-oriented: $W_{\mathrm{pre}} = 1.0$.
3. Recall-oriented: $W_{\mathrm{rec}} = 1.0$.
4. $F_1$-oriented: $W_{\mathrm{f_1}} = 1.0$.
5. Time-oriented: $W_{\mathrm{time}} = 1.0$.

\subsection{Baseline Methods} \label{sec4.5}
We use the following baseline methods:
For evaluation item (i), we use randomly generated rankings as a baseline and average the metrics across 10 random seeds.
For evaluation item (ii), we use five baseline settings. ``Default Setting'' uses LangChain's default configuration, fixed to ``NaiveRAG, stuff, 250, 50, llama3.1:8b, nomic-embed-text'' for all queries. ``Most Frequent Setting'' uses the configuration that appeared most frequently in the empirical $F_1$ and latency evaluations (``NaiveRAG, stuff, 500, 50, llama3.1:8b, nomic-embed-text''). ``Random Selection'' samples a configuration uniformly from the 96-configuration space; its results are averaged over 7,680 trials (80 test queries $\times$ 96 configurations). ``Adaptive-RAG''~\cite{Jeong24} routes queries among No-Retrieval, Single-Step (Default Setting), and Multi-Step (GraphRAG) using a three-class classifier. Finally, ``Oracle'' selects the ex-post optimal configuration from candidate set $\mathcal{C}$ for each query, providing a theoretical upper bound on performance.

\subsection{Evaluation Metrics} \label{sec4.6}
We evaluate RMs using NDCG@K~\cite{JarvelinKalervo02} and Hit@K~\footnote{\url{https://docs.shaped.ai/docs/metrics/hr/}}. NDCG@K measures ranking quality by assigning greater weight to higher-ranked candidates. Hit@K measures the proportion of queries for which the oracle configuration appears among the top $K$ candidates.

To evaluate final RAG response quality, we measure semantic similarity using BERTScore~\cite{zhang20} precision, recall, and $F_1$ score. We also report Exact Match (EM), token-level $F_1$ score, and execution time from query submission to response generation.
To test for statistically significant differences in response accuracy, measured by $F_1$ score, we conduct Wilcoxon signed-rank tests on paired results against key baselines, particularly Adaptive-RAG.

\subsection{Implementation Environment} \label{sec4.7}
We implemented NaiveRAG using LangChain~\footnote{\url{https://www.langchain.com/}} and the graph-based RAG systems using the public LightRAG and GraphRAG repositories on GitHub~\footnote{\url{https://github.com/HKUDS/LightRAG/}}~\footnote{\url{https://github.com/microsoft/graphrag/}}.
For feature extraction, we used spaCy~\footnote{\url{https://spacy.io/}} for token counting, dependency parse-tree depth, and named-entity extraction, and YAKE~\footnote{\url{http://yake.inesctec.pt/}} for keyword extraction.
We ran the experiments on an NVIDIA A100 GPU and used fixed random seeds for reproducibility.
The source code for RM training and inference is publicly available on GitHub (\url{https://anonymous.4open.science/r/Adaptive-RAG-Configuration-Selection-5E65/}).
